\documentclass[11pt]{article}
\usepackage[margin=1in]{geometry}
\usepackage{tabularx,booktabs,array}
\newcolumntype{Y}{>{\raggedright\arraybackslash}X}
\usepackage{amsmath,amssymb}
\usepackage{listings}
\usepackage{xcolor}
\usepackage[hyphens]{url}
\usepackage[htt]{hyphenat}
\usepackage{hyperref}

% Safety net: render the few Greek/accented glyphs that appear inside code
% identifiers (e.g. dξdx) when compiling with pdflatex + utf8.
\usepackage{newunicodechar}
\newunicodechar{ξ}{\ensuremath{\xi}}
\newunicodechar{η}{\ensuremath{\eta}}
\newunicodechar{â}{\ensuremath{\hat{a}}}
\newunicodechar{λ}{\ensuremath{\lambda}}

% Allow long \texttt tokens (with underscores, dots, !) to break across lines
\sloppy
\emergencystretch=3em
\hyphenpenalty=1000

% \code{...} : like \texttt but breakable after every "\_" so that long
% snake_case identifiers are wrapped instead of running into the margin.
\makeatletter
\newcommand{\code}[1]{{\def\_{\char`\_\discretionary{}{}{}}\texttt{#1}}}
\makeatother

\lstset{
    basicstyle=\ttfamily\footnotesize,
    breaklines=true,
    frame=single,
    showstringspaces=false
}

\newcommand{\ahat}{\hat{a}}
\newcommand{\Jk}{J_K}

\title{Mesh-induced reference diffusivity $\ahat$ for Element Learning\\
\large Theory, implementation, and verification}
\author{}
\date{}

\begin{document}
\maketitle

This note documents the \textbf{non-constant diffusivity} capability added to
the Element Learning (EL) solver: what the mesh-induced reference diffusivity
$\ahat$ is, how it is generated and consumed in the code, and how each piece was
verified.

It accompanies the code in
\begin{itemize}
  \item \code{src/kernel/EL\_nonconstant\_diffusivity.jl} --- the $\ahat$ machinery,
  \item \code{src/kernel/elementLearningStructs.jl} --- sampling/inference wiring,
  \item \code{src/io/write\_output.jl} --- VTU visualisation,
  \item \code{problems/Elliptic/elementLearning\_hole\_nonconstant/} --- the geometry test,
  \item \code{problems/Elliptic/case1\_nonconstant/} --- the physical-$a(x,y)$ reference test,
  \item \code{tools/EL\_training/train\_rfrc.py} --- the matching RF-RC trainer.
\end{itemize}

\section{What $\ahat$ is}
\label{sec:what}

On an element $K$ the local stiffness entry of the variable-coefficient operator
$-\nabla\!\cdot(a\,\nabla u)$ is pulled back to the reference element
$[-1,1]^d$:
\begin{align}
(A_{v^o,v})_{ij}
  &= \int_K a(x)\,\nabla_x \phi_i \cdot \nabla_x \phi_j \nonumber\\
  &= \int_{[-1,1]^d} (\ahat\,\nabla_\xi \phi_i)\cdot(\nabla_\xi \phi_j),
\end{align}
with the \textbf{reference diffusivity}
\begin{equation}
\ahat(\xi) \;=\; \frac{|K|}{|[-1,1]^d|}\; a\; \Jk^{-1}\,\Jk^{-\mathsf{T}}
\qquad\text{(a $2\times2$ SPD tensor in 2-D)},
\label{eq:ahat}
\end{equation}
where $\Jk = \partial x/\partial \xi$ is the element Jacobian. $\ahat$ fuses two
ingredients:
\begin{enumerate}
  \item the \textbf{physical} scalar diffusivity $a$ of the PDE, and
  \item the \textbf{geometry} of the element through $\Jk$.
\end{enumerate}
The EL surrogate learns the map
$\ahat \mapsto T^{ie} = (A_{v^o,v^o})^{-1} A_{v^o,v^b}$, i.e.\ the local
interior-elimination operator, replacing an explicit per-element matrix inverse.

\subsection{Two ways $\ahat$ becomes non-constant}

\begin{center}
\begin{tabularx}{\linewidth}{@{}lccY@{}}
\toprule
Source of non-constancy & Physical $a$ & Geometry $\Jk$ & Test case \\
\midrule
\textbf{Physical} ($a(x,y)$) & varies & uniform (structured mesh) & \code{case1\_nonconstant} \\
\textbf{Mesh-induced}        & $a=1$  & varies (non-affine: general quads / curved) & \code{elementLearning\_hole\_nonconstant} \\
\bottomrule
\end{tabularx}
\end{center}

This document focuses on the \textbf{mesh-induced} case, which is the one the EL
notes prescribe (Option~1: \emph{``use an unstructured mesh, which supplies a set
of random element shapes''}). With constant physical $a=1$, the only non-trivial
part of $\ahat$ is the Jacobian --- exactly the \emph{``simplest case''} of the
notes.

\subsection{Constant within an element, or varying? (affine vs.\ bilinear)}
$\ahat$ is constant \emph{within} an element only when the element map is
\textbf{affine}, which for a quadrilateral means a \textbf{parallelogram}
(opposite sides parallel \emph{and} equal); then $\Jk$ is constant.
``Straight sides'' is \emph{not} sufficient: a general straight-sided quad
(trapezoid, kite, or the sheared/fanned quads that conform to a curved boundary
such as the hole) has a \textbf{bilinear} map
$x(\xi,\eta)=\sum_i N_i(\xi,\eta)\,P_i$, so
$\partial N_i/\partial\xi \propto (1\pm\eta)$ and $\Jk=\partial x/\partial\xi$
varies \textbf{linearly inside the element}. Hence $\ahat(\xi)$ --- and the VTU
\code{ahat\_eff}, \code{ahat\_aniso}, \code{ahat\_11/12/22} fields --- \textbf{vary
within each non-parallelogram element}, peaking at the most distorted corners.
The within-element gradient seen in the plots is therefore expected and correct;
only true parallelograms show a single value per element. (Curved/high-order
geometry is just a further source of the same $\xi$-dependence.)

\medskip\noindent\textbf{Sampling consistency.} The default sampler
\code{synthesize\_random\_ahat} draws random \textbf{affine} Jacobians
(parallelograms $\Rightarrow$ $\ahat$ constant per element). To train a surrogate
that matches \emph{general} (bilinear) mesh quads, generate within-element-varying
samples --- \code{:lEL\_xidependent => true}
(\code{synthesize\_random\_ahat\_field}) or a dedicated bilinear-quad sampler.

\subsection{A useful identity (area normalisation)}
\label{sec:detone}

For $a=1$,
\begin{equation}
\det(\ahat)
  = \Big(\tfrac{|K|}{|\mathrm{ref}|}\Big)^{2}\det(\Jk^{-1}\Jk^{-\mathsf{T}})
  = \det(\Jk)^{2}\,\big(1/\det(\Jk)\big)^{2}
  = 1 .
\end{equation}
So \textbf{$\det(\ahat)=1$ identically}: the mesh-induced $\ahat$ is
\emph{area-normalised} and its variation is \textbf{purely anisotropic}. Two
scalar invariants then fully describe it (because
$\lambda_{\max}\lambda_{\min}=\det=1 \Rightarrow \lambda_{\min}=1/\lambda_{\max}$):
\begin{align}
\text{\texttt{ahat\_eff}}   &= \tfrac{1}{2}\,\mathrm{tr}(\ahat)
   = \tfrac{1}{2}(\lambda_{\max}+\lambda_{\min}) \ge 1
   &&(=1 \Leftrightarrow \text{undistorted element}),\\
\text{\texttt{ahat\_aniso}} &= \lambda_{\max}/\lambda_{\min}
   &&(=1 \Leftrightarrow \text{isotropic}).
\end{align}
These are among the fields written to the VTU (Section~\ref{sec:vtu}).

\section{Implementation}

\subsection{Feature representation}
The NN input feature is the per-node $\ahat$ tensor, stored as its three unique
SPD entries $(\ahat_{11},\ahat_{12},\ahat_{22})$ at each of the $(k+1)^2$
reference nodes:
\begin{align*}
\text{feature length} &= 3\,(k+1)^2 \quad\text{(was $(k+1)^2$ for the constant-scalar case)},\\
\text{layout (conn node $m$)} &= [\,\ahat_{11}^{(m)},\ \ahat_{12}^{(m)},\ \ahat_{22}^{(m)}\,].
\end{align*}
The NN output is the flattened $T^{ie}$ of size $n_{vo}\cdot n_{vb}$
($n_{vo}=(k-1)^2$ interior nodes, $n_{vb}=(k+1)^2-(k-1)^2$ boundary nodes),
unchanged.

\subsection{Sampling (training-data generation) --- Option 1}
\code{el\_nonconstant\_sampling!} synthesises random element shapes (random affine
Jacobians: rotation $\times$ anisotropic stretch $\times$ shear) and, per sample,
\begin{enumerate}
  \item forms $\ahat = a\,\det(J)\,(J^{\mathsf{T}}J)^{-1}$ (constant per element
        for an affine/parallelogram shape, or varying within the element ---
        general bilinear quad / curved --- via
        \code{synthesize\_random\_ahat\_field}),
  \item assembles the local reference stiffness
        (\code{el\_assemble\_local\_stiffness} via the precomputed components
        \code{K11,K12,K22}, or \code{el\_assemble\_local\_stiffness\_field} for the
        $\xi$-dependent case),
  \item condenses $T^{ie} = (A_{v^o v^o})^{-1} A_{v^o v^b}$,
  \item writes $(\text{feature},\ \mathrm{vec}(T^{ie}))$ in \textbf{\code{mesh.conn} node order}.
\end{enumerate}

\subsection{Inference}
\code{el\_avisc\_nonconstant!} builds the per-element feature from the \textbf{real
mesh metrics} $\ahat = J_e\, a\, \Jk^{-1}\Jk^{-\mathsf{T}}$ (with
$\Jk^{-1} = [\,d\xi dx\ \ d\xi dy;\ d\eta dx\ \ d\eta dy\,]$,
$a=$\,\code{el\_diffusivity}$(x,y)$, default $1$), in the same \code{mesh.conn} order.
Because sampling and inference share that single ordering,
\code{elementLearning\_infer!} is reused \textbf{unchanged} --- only the feature
differs.

Gated by \code{inputs[:lEL\_nonconstant]} (default \code{false}), so the existing
constant-amplitude pipeline and its trained model are untouched.

\subsection{Operator for the physical reference test}
For \code{case1\_nonconstant}, \code{build\_laplace\_matrix(NSD\_2D; afun)} multiplies
the integrand by $a(x,y)$ at each node so the discrete operator is genuinely
$-\nabla\!\cdot(a\nabla u)$; \code{afun === nothing} reproduces the plain
Laplacian exactly.

\subsection{VTU visualisation}
\label{sec:vtu}
\code{write\_vtk(NSD\_2D)} writes the geometry-induced $\ahat$ to the VTU when
\code{inputs[:lEL\_nonconstant]} is set and \code{metrics} are supplied.
\textbf{Three output formats} are available via \code{inputs[:ahat\_output]}
(default \code{:cell}):

\begin{center}
\begin{tabularx}{\linewidth}{@{}lllY@{}}
\toprule
\code{inputs[:ahat\_output]} & VTU fields & Location & Description \\
\midrule
\code{:cell} (default) & \code{ahat\_eff}, \code{ahat\_aniso} & cell data &
  invariants, element-averaged (piecewise-constant per element) \\
\code{:nodal} & \code{ahat\_eff}, \code{ahat\_aniso} & point data &
  same invariants, \textbf{smoothed nodal} field ($\ahat$ averaged over the
  elements sharing each node, then invariants taken) \\
\code{:tensor} & \code{ahat\_11}, \code{ahat\_12}, \code{ahat\_22} & point data &
  the full smoothed nodal \textbf{tensor components} \\
\bottomrule
\end{tabularx}
\end{center}

\noindent
with $\text{\texttt{ahat\_eff}}=\tfrac{1}{2}\mathrm{tr}(\ahat)=\tfrac{1}{2}(\lambda_{\max}+\lambda_{\min})\ge1$
($1\Leftrightarrow$ undistorted) and
$\text{\texttt{ahat\_aniso}}=\lambda_{\max}/\lambda_{\min}$
($1\Leftrightarrow$ isotropic). The smoothing averages the three $\ahat$
components over the elements incident to each node (count-weighted).

For the \emph{physical} case, \code{case1\_nonconstant} instead writes the nodal
scalar field \code{diffusivity}~$=a(x,y)$ (gated on \code{inputs[:diffusivity]}).

\section{Verification}
\label{sec:verify}

The full JEXPRESSO solve (GMSH mesh $+$ trained ONNX model) was \textbf{not} run
in this environment, so verification targets the mathematics and the data flow
with standalone Julia checks. All tolerances below are machine precision
($\sim 10^{-16}$--$10^{-10}$) unless noted; all checks \textbf{passed}.

\subsection{Manufactured solution (physical reference, \code{case1\_nonconstant})}
For $a(x,y)=1+0.2x$, $u_{\mathrm{ex}}=\sin(x)\cos(y)$, the analytic source
\begin{equation*}
f = -A_1 k_x \cos(k_x x)\cos(k_y y)
    + (A_0 + A_1 x)(k_x^2+k_y^2)\sin(k_x x)\cos(k_y y)
\end{equation*}
was checked against a central finite-difference evaluation of
$-\nabla\!\cdot(a\nabla u_{\mathrm{ex}})$:
\begin{equation*}
\max\big|\,{-}\nabla\!\cdot(a\nabla u_{\mathrm{ex}})_{\mathrm{FD}} - f\,\big|
   \approx 2.6\times10^{-6}\quad(\text{FD truncation; PASS}),
\end{equation*}
with $a\in[1.2,2.2]>0$ on the domain.

\subsection{Diffusivity-weighted operator (\code{build\_laplace\_matrix})}
Replicating the assembly on a reference element:
\begin{itemize}
  \item \code{afun === nothing} reproduces the pure reference Laplacian (matches
        the EL module's \code{K11}\,+\,\code{K22}) $\Rightarrow$ \textbf{existing
        problems unaffected};
  \item constant $a=c$ scales the operator by exactly $c$;
  \item non-constant $a(x,y)$ keeps the element matrix \textbf{symmetric, PSD,
        with the constant nullspace}.
\end{itemize}

\subsection{Local $\ahat$ machinery (\code{EL\_nonconstant\_diffusivity.jl})}
\begin{itemize}
  \item the reference Laplacian ($\ahat=I$) is symmetric, PSD, constant nullspace;
  \item the precomputed $K$-component assembly equals the full $O(N^6)$
        quadrature assembly for constant $\ahat$;
  \item synthesised $\ahat$ is \textbf{SPD on every draw} (2000 samples);
  \item $T^{ie}$ reproduces the local Dirichlet interior recovery
        $u_{v^o} = -T^{ie} u_{v^b}$ vs the direct block solve;
  \item the metrics-based $\ahat$ (inference) equals the synthesised $\ahat$
        (sampling) for the \textbf{same Jacobian} --- i.e.\ \textbf{sampling
        $\leftrightarrow$ inference consistency};
  \item feature/output dimensions are $3(k+1)^2$ and $n_{vo}\cdot n_{vb}$
        (e.g.\ $nop=6 \to 147$ and $25\times24$).
\end{itemize}

\subsection{Inference-wiring node ordering}
\begin{itemize}
  \item the \code{mesh.conn} $\leftrightarrow$ tensor $(i,j)$ map round-trips on
        the element layout;
  \item the conn-ordered $T^{ie}$ reproduces the direct interior recovery,
        confirming \code{elementLearning\_infer!} can be reused for the
        non-constant feature;
  \item $\xi$-dependent (curved-element) $\ahat$ is SPD and genuinely
        node-varying, and its assembled stiffness is symmetric/PSD with the
        constant nullspace.
\end{itemize}

\subsection{VTU $\ahat$ formats (\code{write\_vtk})}
\begin{itemize}
  \item \textbf{undistorted} element $J=s\,I$: $\text{\texttt{ahat\_eff}}=1$,
        $\text{\texttt{ahat\_aniso}}=1$, $\det(\ahat)=1$ at \textbf{all scales} $s$
        (scale-invariant, as expected);
  \item \textbf{arbitrary distorted} $J$ (rotation/stretch/shear sweep):
        $\det(\ahat)=1$ always, $\ahat$ SPD, $\text{\texttt{ahat\_eff}}\ge1$,
        $\text{\texttt{ahat\_aniso}}\ge1$;
  \item \code{:cell} gives distinct per-element values; the \textbf{smoothed
        nodal} tensor (\code{:tensor}) at a shared edge equals the mean of the
        adjacent elements' $\ahat$; \code{:nodal} invariants from the smoothed
        tensor satisfy $\text{\texttt{ahat\_eff}}\ge1$;
  \item the VTU invariants are computed from the \textbf{same} $\ahat$ entries as
        \code{el\_ahat\_nodes\_from\_metrics}, so \emph{what you see equals what the
        surrogate consumes}.
\end{itemize}

\subsection{Trainer}
\code{tools/EL\_training/train\_rfrc.py} auto-sizes \code{N\_in} from the CSV (so
$3(k+1)^2$ is picked up automatically) and standardises the heterogeneous
$\ahat$ feature, baking the standardisation into the exported ONNX so the Julia
side keeps feeding the \textbf{raw} $\ahat$ feature. The
standardise\,$\to$\,bake-in round-trip (\code{export(raw)} \texttt{==}
\code{train(standardised)}) was verified.

\section{How to run}

\subsection{Geometry-induced $\ahat$ test (the subject of this note)}
\begin{lstlisting}
julia --project=.
julia> push!(empty!(ARGS), "Elliptic", "elementLearning_hole_nonconstant");
julia> include("./src/Jexpresso.jl")
\end{lstlisting}
\begin{itemize}
  \item \textbf{Verify the physical problem and see $\ahat$ now} (no model
        needed): comment \code{:lelementLearning} in \code{user\_inputs.jl}. The
        direct solver solves $-\Delta u = f$ on the hole, prints the MMS
        $\lVert e\rVert_{L^2}$, and writes \code{ahat\_eff}/\code{ahat\_aniso} (or
        the tensor components, per \code{inputs[:ahat\_output]}) to
        \code{output-nonconstant/iter\_1.pvtu}. The anisotropy concentrates around
        the circular hole, where elements are most distorted.
  \item \textbf{EL surrogate inference}: keep \code{:lelementLearning => true}
        with a model trained on the $3(k+1)^2$ $\ahat$ feature (sample with
        \code{:lEL\_Sample}, train with \code{tools/EL\_training/train\_rfrc.py}).
\end{itemize}

\subsection{Physical $a(x,y)$ reference test}
\begin{lstlisting}
julia> push!(empty!(ARGS), "Elliptic", "case1_nonconstant");
julia> include("./src/Jexpresso.jl")
\end{lstlisting}
Direct solve of $-\nabla\!\cdot(a\nabla u)=f$ with $a=1+0.2x$; prints the MMS
$\lVert e\rVert_{L^2}$ and writes the nodal \code{diffusivity} field to the VTU.

\section{Known limitations / follow-ups}
\begin{itemize}
  \item The EL \textbf{inference} path needs a model retrained on the
        $3(k+1)^2$ $\ahat$ feature; running with the old $(k+1)^2$ model $+$
        \code{:lEL\_nonconstant} fails on a dimension mismatch (expected).
  \item For a \emph{physical} non-constant $a(x,y)$ solved \textbf{through} the
        EL surrogate, the global skeleton matrix $L$ should also be assembled
        with that $a$ (use \code{inputs[:diffusivity]}); the mesh-induced
        ($a=1$) case is already self-consistent because $L$ is the standard
        Laplacian.
  \item VTU $\ahat$ fields and the diffusivity operator are \textbf{2-D}
        (\code{NSD\_2D}), matching the current EL scope.
  \item Three $\ahat$ VTU formats are provided (\code{inputs[:ahat\_output]} =
        \code{:cell} / \code{:nodal} / \code{:tensor}, Section~\ref{sec:vtu}).
        The nodal smoothing is a count-weighted average over incident elements;
        across MPI partitions it is a per-partition average (a visualisation
        field, not used in the solve).
\end{itemize}

\end{document}
